% MARQ Network Architecture - TikZ Diagram
% DeepNash-style visualization for academic papers
% Compile with: pdflatex marq_network_diagram.tex

\documentclass[tikz,border=10pt]{standalone}
\usepackage{tikz}
\usetikzlibrary{positioning,shapes.geometric,arrows.meta,calc,fit,backgrounds,decorations.pathreplacing}
\usepackage{xcolor}

% Define MARQ color palette
\definecolor{convblue}{RGB}{52,152,219}
\definecolor{convblue2}{RGB}{41,128,185}
\definecolor{attnorange}{RGB}{230,126,34}
\definecolor{pbspurple}{RGB}{155,89,182}
\definecolor{valugreen}{RGB}{46,204,113}
\definecolor{advred}{RGB}{231,76,60}
\definecolor{noisyyellow}{RGB}{241,196,15}
\definecolor{darkbg}{RGB}{44,62,80}
\definecolor{lightgray}{RGB}{189,195,199}

\begin{document}
\begin{tikzpicture}[
    % Styles
    convlayer/.style={
        rectangle,
        draw=convblue2,
        fill=convblue,
        minimum width=1.2cm,
        minimum height=2cm,
        text=white,
        font=\footnotesize\bfseries,
        rounded corners=2pt
    },
    convlayer3d/.style={
        convlayer,
        path picture={
            \fill[convblue2] (path picture bounding box.north east) -- 
                ++(0.15,-0.15) -- ++(0,-1.7) -- (path picture bounding box.south east) -- cycle;
            \fill[convblue!70!black] (path picture bounding box.north west) -- 
                (path picture bounding box.north east) -- ++(0.15,-0.15) -- ++(-1.2,0) -- cycle;
        }
    },
    noisylayer/.style={
        rectangle,
        draw=noisyyellow!80!black,
        fill=noisyyellow,
        minimum width=0.8cm,
        minimum height=1.5cm,
        text=darkbg,
        font=\tiny\bfseries,
        rounded corners=2pt,
        decorate,
        decoration={random steps, segment length=2pt, amplitude=1pt}
    },
    valstream/.style={
        rectangle,
        draw=valugreen!80!black,
        fill=valugreen,
        minimum width=1.5cm,
        minimum height=1.2cm,
        text=white,
        font=\footnotesize\bfseries,
        rounded corners=3pt
    },
    advstream/.style={
        rectangle,
        draw=advred!80!black,
        fill=advred,
        minimum width=1.5cm,
        minimum height=1.2cm,
        text=white,
        font=\footnotesize\bfseries,
        rounded corners=3pt
    },
    attention/.style={
        ellipse,
        draw=attnorange!80!black,
        fill=attnorange,
        minimum width=2cm,
        minimum height=1cm,
        text=white,
        font=\footnotesize\bfseries
    },
    pbsbox/.style={
        rectangle,
        draw=pbspurple!80!black,
        fill=pbspurple,
        minimum width=2.5cm,
        minimum height=1.5cm,
        text=white,
        font=\footnotesize\bfseries,
        rounded corners=3pt
    },
    inputbox/.style={
        rectangle,
        draw=lightgray,
        fill=darkbg,
        minimum width=1.8cm,
        minimum height=1.8cm,
        text=white,
        font=\footnotesize
    },
    outputbox/.style={
        rectangle,
        draw=advred,
        fill=advred!90!black,
        minimum width=2cm,
        minimum height=0.8cm,
        text=white,
        font=\footnotesize\bfseries,
        rounded corners=5pt
    },
    arrow/.style={
        ->,
        >=Stealth,
        thick,
        draw=lightgray
    },
    label/.style={
        font=\scriptsize,
        text=darkbg
    }
]

% ============ INPUT SECTION ============
\node[inputbox, label=below:{\footnotesize Board State}] (board) at (0,0) {
    \begin{tabular}{c}
    \textbf{10×10}\\
    Stratego
    \end{tabular}
};

\node[inputbox, below=0.5cm of board, label=below:{\footnotesize PBS Beliefs}] (pbs) {
    \begin{tabular}{c}
    \textbf{12 ch}\\
    P(type)
    \end{tabular}
};

% Combine into 27-channel input
\node[rectangle, draw=convblue, fill=convblue!20, minimum width=1.5cm, minimum height=4cm, 
      right=1cm of $(board)!0.5!(pbs)$, rounded corners=3pt, label=above:{\footnotesize 27 Channels}] (input27) {};
\node[font=\tiny, text=convblue2] at (input27) {
    \begin{tabular}{c}
    PBS (12)\\
    +\\
    Own (12)\\
    +\\
    Masks (3)
    \end{tabular}
};

\draw[arrow] (board.east) -- (input27.west |- board);
\draw[arrow] (pbs.east) -- (input27.west |- pbs);

% ============ CONVOLUTIONAL ENCODER ============
\node[convlayer3d, right=1.2cm of input27, minimum height=3cm, label=above:{\footnotesize Conv 32}] (conv1) {32};
\node[convlayer3d, right=0.4cm of conv1, minimum height=2.5cm, minimum width=1cm, label=above:{\footnotesize Conv 64}] (conv2) {64};
\node[convlayer3d, right=0.4cm of conv2, minimum height=2cm, minimum width=0.8cm, label=above:{\footnotesize Conv 64}] (conv3) {64};

\draw[arrow] (input27.east) -- (conv1.west);
\draw[arrow] (conv1.east) -- (conv2.west);
\draw[arrow] (conv2.east) -- (conv3.west);

% Flatten indicator
\node[rectangle, draw=lightgray, dashed, minimum width=0.3cm, minimum height=1.5cm, 
      right=0.3cm of conv3] (flatten) {};
\node[font=\tiny, rotate=90] at (flatten) {Flatten};
\draw[arrow] (conv3.east) -- (flatten.west);

% ============ NOISY LAYERS ============
\node[noisylayer, right=0.5cm of flatten, label=above:{\tiny Noisy FC}] (noisy1) {512};
\node[noisylayer, right=0.3cm of noisy1, label=above:{\tiny Noisy FC}] (noisy2) {512};

\draw[arrow] (flatten.east) -- (noisy1.west);
\draw[arrow] (noisy1.east) -- (noisy2.west);

% ============ DUELING ARCHITECTURE ============
% Value stream
\node[valstream, above right=0.5cm and 1cm of noisy2] (value) {
    \begin{tabular}{c}
    Value\\
    $V(s)$
    \end{tabular}
};

% Advantage stream
\node[advstream, below right=0.5cm and 1cm of noisy2] (advantage) {
    \begin{tabular}{c}
    Advantage\\
    $A(s,a)$
    \end{tabular}
};

\draw[arrow] (noisy2.east) -- ++(0.3,0) |- (value.west);
\draw[arrow] (noisy2.east) -- ++(0.3,0) |- (advantage.west);

% Combine node
\node[circle, draw=pbspurple, fill=pbspurple!30, minimum size=0.8cm, 
      right=1.5cm of $(value)!0.5!(advantage)$, font=\footnotesize] (combine) {$+$};

\draw[arrow] (value.east) -- ++(0.3,0) |- (combine.west);
\draw[arrow] (advantage.east) -- ++(0.3,0) |- (combine.west);

% Q-value output
\node[rectangle, draw=pbspurple, fill=pbspurple!80, minimum width=2cm, minimum height=1cm,
      right=0.8cm of combine, text=white, font=\footnotesize\bfseries, rounded corners=3pt] (qvalue) {
    \begin{tabular}{c}
    $Q(s,a)$\\
    {\tiny 51 atoms}
    \end{tabular}
};

\draw[arrow] (combine.east) -- (qvalue.west);

% Action output
\node[outputbox, right=0.8cm of qvalue] (action) {Action $a^*$};
\draw[arrow, draw=advred] (qvalue.east) -- (action.west);

% ============ AAREN MODULE (BELOW) ============
\node[attention, below=2.5cm of conv2] (aaren) {AAREN};

\node[rectangle, draw=attnorange!50, fill=attnorange!20, minimum width=3cm, minimum height=0.6cm,
      left=1cm of aaren, font=\tiny, rounded corners=2pt] (history) {
    Observation History $\{o_1, \ldots, o_t\}$
};

\node[pbsbox, right=1cm of aaren] (pbsupdate) {
    \begin{tabular}{c}
    PBS Update\\
    $\mathcal{B}_t$
    \end{tabular}
};

\draw[arrow, draw=attnorange] (history.east) -- (aaren.west);
\draw[arrow, draw=attnorange] (aaren.east) -- (pbsupdate.west);
\draw[arrow, draw=pbspurple, dashed] (pbsupdate.north) -- ++(0, 1.5) -| (input27.south);

% ============ LABELS AND ANNOTATIONS ============
% Main section labels
\node[font=\small\bfseries, text=darkbg] at ($(conv1)!0.5!(conv3) + (0, 2)$) {Convolutional Encoder};
\node[font=\small\bfseries, text=noisyyellow!80!black] at ($(noisy1)!0.5!(noisy2) + (0, 1.5)$) {Noisy Layers};
\node[font=\small\bfseries, text=pbspurple] at ($(value)!0.5!(advantage) + (3, 0)$) {Dueling};

% Decorative brace for dueling architecture
\draw[decorate, decoration={brace, amplitude=5pt, raise=5pt}, draw=pbspurple!70]
    (value.north east) -- (advantage.south east) 
    node[midway, right=10pt, font=\tiny, text=pbspurple] {$Q = V + (A - \bar{A})$};

% C51 label
\node[font=\tiny, text=pbspurple, below=0.1cm of qvalue] {Categorical DQN};

% Title
\node[font=\Large\bfseries, text=darkbg] at ($(conv2) + (3, 4)$) {MARQ: Rainbow DQN Architecture};

\end{tikzpicture}
\end{document}
